\begin{problem}{개성있는 구간}{standard input}{standard output}{1 second}{256 megabytes}

정수 $N$이 주어진다.


다음 조건을 만족하는 길이 $N$의 정수 수열 $A = [A_1, A_2, \dots, A_N]$을 구하여라

\begin{itemize}
    \item 모든 연속 부분수열의 합이 서로 다르다.  
    즉, 서로 다른 두 구간 $[L_1, R_1]$, $[L_2, R_2]$에 대해
    $$
    \sum_{i=L_1}^{R_1} A_i \ne \sum_{i=L_2}^{R_2} A_i
    $$
    가 항상 성립해야 한다.
    \item $1 \le i \le N$인 모든 $i$에 대해 $1 \le A_i \le 6N$ 이다.
\end{itemize}

여기서 구간 $[L, R]$은 $1 \le L \le R \le N$을 만족하는 연속 부분수열을 의미한다.

조건을 만족하는 수열이 여러 개 존재할 수 있다. 그 중 아무거나 하나 출력하면 된다.

\InputFile
첫째 줄에 정수 $N$ $(1 \le N \le 8192)$이 주어진다.

\OutputFile
첫째 줄에 정수 $N$개 $A_1, A_2, \dots, A_N$을 공백으로 구분하여 출력하라.
$(1 \le A_i \le 6N)$


\Example

\begin{example}
\exmpfile{example.01}{example.01.a}%
\end{example}

\end{problem}

